\chapter{Methodology}
\label{ch:methodology}

This chapter describes the methodology used to conduct the State-of-the-Practice study of MPS software and addresses RQ1 (candidate identification) by documenting how the candidate set was assembled. The methodology follows \citet{Smith2021} and adapts it to the MPS domain. Section~\ref{sec:domain-selection} explains how the domain was chosen, and Section~\ref{sec:domain-expert-interaction} describes the role of the domain expert. Section~\ref{sec:candidate-software-discovery} covers the three-source candidate discovery process; Section~\ref{sec:mention-count-ranking} introduces the mention-count ranking that orders those candidates; and Section~\ref{sec:filtering-criteria} sets out the filtering criteria applied to the ranked list. The resulting selected set of 28 packages appears in Section~\ref{sec:final-software-set}. Section~\ref{sec:measurement-procedure} then walks through the measurement procedure using the Domain~X template, Section~\ref{sec:commonality-analysis-procedure} outlines the commonality analysis, and Section~\ref{sec:interview-component} describes the developer interview component.

\section{Domain Selection}
\label{sec:domain-selection}

The first step of the Domain~X methodology is to identify a suitable research software domain. To be applicable for the methodology, the chosen domain must meet four criteria~\citep{Smith2021}: (i)~it must have a well-defined and stable theoretical underpinning, ideally supported by standard textbooks; (ii)~there must be a community of researchers and practitioners studying the domain; (iii)~the software packages in the domain must include open-source options; and (iv)~a preliminary search should suggest that there are close to 30 candidate software packages that qualify as research software.

The selected domain is \emph{MPS (motion planning and simulation)}. Motion planning and simulation are established subfields of robotics with well-defined theoretical foundations~\citep{lavalle2006planning,choset2005principles}. There is an active international research community, regular conferences (such as ICRA, IROS, and RSS), and a substantial body of published work. A preliminary search confirmed the existence of numerous open-source software packages spanning simulation environments, physics engines, dynamics libraries, motion planning libraries, and robotics middleware.

The domain boundary follows the MPS scope defined in Chapter~\ref{ch:background}, Section~\ref{sec:robotics-software-background}, rather than a simulator-only scope.

Earlier project discussions considered a narrower domain focused on inverse kinematics software. After consultation with the supervisor and domain expert, the scope was broadened to the MPS domain, which provides a richer set of candidate packages and better reflects the interconnected nature of robotics software ecosystems.

\section{Interaction with the Domain Expert}
\label{sec:domain-expert-interaction}

The Domain~X methodology emphasizes that a domain expert is an essential member of the assessment team. Pitfalls exist if non-experts attempt to produce an authoritative list of software or definitively rank packages without domain knowledge. Non-experts must rely on information available online, which has two known drawbacks: (i)~online resources may contain false or inaccurate information; and (ii)~online resources may omit relevant information that is so ingrained with experts that nobody thinks to record it explicitly~\citep{Smith2021}.

Domain experts may be recruited from academia or industry. The only requirements are knowledge of the domain and a willingness to be engaged in the assessment process. The domain expert does not need to be a software developer, but they should be a user of domain software. Given that domain experts are likely to be busy people, the methodology is designed to minimize the burden on their time~\citep{Smith2024MI}.

For this report, the domain expert is Dr.~Matthew Giamou of the Department of Computing and Software at McMaster University, whose research includes robotics, motion planning, and state estimation. Dr.~Giamou reviewed the candidate software list and provided feedback on scope and coverage. He confirmed that the list was reasonable and that most of the shortlisted tools were familiar to him. He noted that several packages (ROS~2, MoveIt!, and OMPL) are not full simulators but can fit within a coherent narrative if their roles are explained carefully. He also confirmed that excluding MATLAB/Simulink, Unity, NVIDIA Isaac Sim, and RaiSim on methodological grounds (closed source or non-OSS licensing for core components, as detailed in Section~\ref{subsec:open-source-availability}) was acceptable and would not misrepresent the MPS domain. MRDS and USARSim, which are also excluded, were filtered on independent activity grounds (Section~\ref{subsec:activity-maintenance-status}) rather than as a separate expert decision.

\section{Candidate Software Discovery}
\label{sec:candidate-software-discovery}

Candidate software packages were identified from three types of sources: academic survey papers, community-curated software lists, and LLM-generated software lists. Using multiple source types reduces dependence on any single search method and provides broader coverage of the MPS domain. The collected names were consolidated to merge equivalent entries (for example, ``V-REP'' and ``CoppeliaSim'' refer to the same package), and the consolidated list was used as input to the mention-count ranking described in Section~\ref{sec:mention-count-ranking}.

\subsection{Academic Sources}
\label{subsec:academic-sources}

Academic sources identify software packages that appear in research discussions of robotics simulators, simulation tools, and related challenges. These sources included survey papers and comparative studies, such as \citet{Harris2011}, \citet{collins2021simulators}, \citet{Korber2021}, and \citet{Farley2022}, that discuss packages such as Gazebo, Webots, V-REP/CoppeliaSim, MuJoCo, and PyBullet in the context of robotics research. The full list of academic sources used for mention counting appears in Appendix~\ref{appendix:mention-count-source-list}. Academic sources connect software packages to published research contexts and provide a research-facing view of the MPS domain.

Academic sources have known limitations as the sole basis for candidate discovery. Survey papers tend to emphasize packages relevant to a specific research question, robot type, or time period; some active tools may not appear in any of the surveys consulted, while older tools sometimes remain visible because they were prominent when a survey was written. For that reason, academic sources were combined with community-curated and LLM-assisted sources instead of being used on their own.

\subsection{Community-Curated Sources}
\label{subsec:community-curated-sources}

Community-curated sources broaden the view of robotics software practice beyond what appears in academic surveys. The source catalogue included GitHub Topic pages, ``awesome'' lists (curated collections such as \texttt{awesome-robotics-libraries}, \texttt{awesome-robotic-tooling}, and \texttt{awesome-motion-planning}), best-of lists such as \texttt{best-of-robot-simulators}, and other community-maintained aggregations. Sources covered topics including motion planning, robotics simulation, robotics libraries, ROS~2 ecosystem tools, multibody dynamics simulation, and general robotic tooling. The GitHub Topic pages were manually inspected as dynamic community aggregations, while the static curated-list pages were preserved as copied raw web material where available. Appendix~\ref{appendix:mention-count-source-list} lists the community-curated sources individually.

These sources collect software that practitioners and community contributors consider relevant. They can reveal projects that are underrepresented in academic surveys and can show relationships between simulators, planning libraries, middleware, physics engines, and tools. Community-curated sources also vary in maintenance quality, inclusion criteria, and update frequency. They were therefore used as discovery aids rather than as indicators of software quality.

\subsection{LLM-Assisted Sources}
\label{subsec:llm-assisted-sources}

LLM-generated software lists were used as an auxiliary source for candidate discovery. This modifies the original Domain~X methodology, which prescribes search engine queries and domain expert consultation as the primary discovery methods. The modification is recorded explicitly so that the methodology remains transparent and reproducible in principle.

The reason for including LLM-assisted discovery is practical. Large language models surface software names that appear across a wide range of public information, and they act as a broad aggregator over that data. LLM output can also be incomplete, outdated, or incorrect, so names obtained from LLM-generated lists were consolidated, filtered, and verified against public evidence before any claims in this report were based on them. In this report, LLM-assisted sources contributed to candidate discovery but did not determine the selected set, and they were never treated as authoritative.

This approach was discussed with the supervisor. The use of LLMs is technically a modification of the original methodology, but it is defensible for two reasons: (i)~LLMs were used together with academic and community-curated sources, not as the sole discovery method; and (ii)~an LLM, in this context, is a large-scale aggregator over existing public information, analogous in effect to a broad search over indexed web content.

\section{Mention Count Ranking}
\label{sec:mention-count-ranking}

After consolidating equivalent names across the collected sources, each unique software package was counted once per source to produce a mention count. A package that appeared in five academic papers and three community-curated lists would receive a mention count of eight. The consolidated ranking was recorded during the discovery phase and is a transparent input to the filtering process.

Consolidation produced 295 unique candidate names. The top entries in the consolidated ranking include Gazebo (22 mentions), Webots (16), Bullet/PyBullet (16), CoppeliaSim/V-REP (13), MuJoCo (12), ODE (10), ROS~2 (10), and OpenRAVE (9). These counts help organize the candidate list and make the selection process auditable. Appendix~\ref{appendix:full-candidate-software-list} records every candidate with at least four mentions, together with its mention count and filtering outcome, and Appendix~\ref{appendix:mention-count-source-list} lists the counted sources. The full ranked list of all 295 names, including packages that did not survive the filtering step described in Section~\ref{sec:filtering-criteria}, is preserved in the supplementary project materials that accompany this report (a consolidated ranking spreadsheet and a companion source catalogue).

Mention count is a source-appearance signal and must be interpreted accordingly. It does not measure software quality, user adoption, research impact, maintainability, or the overall suitability of any particular software package. A package may appear frequently because it is historically established, because it is broadly useful across application areas, or simply because the sources collected happened to surface it. A more specialized or newer package may be under-represented even when it is technically strong. The mention count is one input to candidate filtering, not the final ranking of the software.

\section{Filtering Criteria}
\label{sec:filtering-criteria}

The ranked candidate list was filtered to identify packages suitable for repository-based State-of-the-Practice measurement. In the Domain~X methodology, the initial filters are scope, usage, and age, applied in priority order until the list reaches a manageable size of approximately 30 packages~\citep{Smith2021}. In this report, these filters are implemented through three criteria: open-source availability, repository-based measurability, and activity and maintenance status. Both the initial and filtered lists are preserved for traceability.

\subsection{Open Source Availability}
\label{subsec:open-source-availability}

Open-source availability was required because the study depends on public software artifacts for measurement. The Domain~X candidate criteria require that source code be viewable and express a preference for GitHub-style repositories from which empirical measures can be collected~\citep{Smith2021}. Source code, repository history, issue trackers, documentation, and release information provide the evidence base for the measurements reported in Chapter~\ref{ch:measurement-results}.

This criterion led to the exclusion of several prominent robotics tools whose core components are not open source:

\begin{itemize}
    \item \textbf{MATLAB and Simulink} are major platforms in both robotics research and industry, widely used for algorithm development, control design, and simulation. However, MATLAB is proprietary software. Its source code is not publicly available, and its development process is not accessible through public repositories. These properties prevent the kind of repository-based measurement that this report depends on.
    \item \textbf{Unity} is a widely used game engine with robotics simulation capabilities, but its core simulation engine is proprietary.
    \item \textbf{NVIDIA Isaac Sim} builds on NVIDIA Omniverse and provides advanced robotics simulation features, but its core platform components are not open source.
    \item \textbf{RaiSim} is a high-performance physics simulator for robotics, but its source code is not publicly available under an open-source license.
\end{itemize}

The exclusions above are methodological constraints, not judgements about the quality or importance of the excluded tools. Each of these platforms plays a significant role in robotics research and industry, and their exclusion limits the coverage of the study. Chapter~\ref{ch:threats-to-validity} discusses this limitation further.

\subsection{Repository-Based Measurability}
\label{subsec:repository-based-measurability}

Repository-based measurability requires that a package have sufficient public evidence to support the measurement process described in Section~\ref{sec:measurement-procedure}. Relevant evidence includes source repositories, commit history, issue data, documentation, release information, build or installation instructions, test artifacts, and project metadata. Packages without sufficient public evidence cannot be measured consistently and were excluded.

This criterion also governs how multi-repository projects are handled. Some robotics software systems are distributed across multiple repositories or GitHub organizations. In such cases, the repositories most representative of the core software were selected for measurement, and this choice was documented as part of the measurement record. The same principle applies to projects that have changed names, migrated repositories, or spawned successor projects.

\subsection{Activity and Maintenance Status}
\label{subsec:activity-maintenance-status}

Activity and maintenance status were considered because the study focuses on contemporary software practice. In the Domain~X methodology, age is one of the initial filtering criteria, measured by the last date when a change was made, with exceptions for older packages that remain highly recommended and currently in use~\citep{Smith2021}. Tools that are no longer maintained provide weaker evidence about current development practices.

The primary activity indicator used in this report is the Domain~X alive/dead rule: a package is considered alive if it has received commits or version releases within the last 18 months from the measurement date. The following packages were excluded on maintenance grounds:

\begin{itemize}
    \item \textbf{MRDS} (Microsoft Robotics Developer Studio) has not received updates since 2012 and is no longer supported by Microsoft.
    \item \textbf{USARSim} (Unified System for Automation and Robot Simulation) is no longer actively maintained and does not represent current practice.
\end{itemize}

\subsection{Scope and Coverage Filters}
\label{subsec:scope-coverage-filters}

Two algorithm-level reference implementations were excluded to keep the selected packages methodologically coherent rather than for quality reasons. \textbf{GPMP2} (Gaussian Process Motion Planner 2) and \textbf{CHOMP} (Covariant Hamiltonian Optimization for Motion Planning) are optimization-based motion planning algorithms with public reference implementations. They are relevant to motion planning research but are narrower in scope than the general-purpose simulators, physics engines, planning frameworks, and middleware tools that dominate the candidate list. Including them alongside Gazebo, MoveIt!, or ROS~2 would force a comparison between full software stacks and single-algorithm libraries, which would distort the commonality analysis and the cross-package measurements. They were therefore excluded to maintain methodological consistency, not because of any quality concern.

Filtering decisions, including the alive/dead status used in Section~\ref{subsec:activity-maintenance-status}, are based on a measurement-time snapshot of repository activity, last commit dates, and release status. These signals can change after measurement, so the measurement records in Chapter~\ref{ch:measurement-results} include the measurement date.

\section{Final Software Set}
\label{sec:final-software-set}

After ranking by mention count and applying the filtering criteria described above, the selected set consists of 28 software packages. Appendix~\ref{appendix:full-candidate-software-list} records the disposition of the leading candidates, and Appendix~\ref{appendix:excluded-software-packages} details the excluded tools. Figure~\ref{fig:candidate-selection} summarizes the full selection pipeline, from the three discovery sources through the mention-count ranking and the filtering criteria to the final set. Table~\ref{tab:final-software-set} lists the packages alphabetically with their primary roles in the robotics software ecosystem.

\begin{figure}[htbp]
\centering
\begin{tikzpicture}[
    font=\small,
    box/.style={draw, rounded corners, align=center, inner sep=5pt, minimum height=0.85cm},
    src/.style={box, fill=blue!5, text width=2.9cm},
    proc/.style={box, fill=black!4, text width=8.4cm},
    final/.style={box, fill=blue!12, text width=5.2cm, font=\small\bfseries},
    excl/.style={box, draw=orange!60!black, fill=orange!6, align=left, text width=3.9cm, font=\footnotesize},
    arr/.style={-{Latex[length=2.5mm]}, thick},
    node distance=9mm and 7mm,
]
\node[src] (s1) {Academic\\survey papers};
\node[src, right=of s1] (s2) {Community-curated\\lists};
\node[src, right=of s2] (s3) {LLM-assisted\\lists};
\node[proc, below=12mm of s2] (cons) {Consolidate equivalent names, then rank by mention count\\[2pt]\footnotesize e.g.\ Gazebo~22, Webots~16, Bullet/PyBullet~16,\\\footnotesize CoppeliaSim/V-REP~13, MuJoCo~12, ODE~10, ROS~2~10, OpenRAVE~9};
\node[proc, below=of cons] (filt) {Filtering criteria (applied in priority order):\\[2pt]open-source availability $\cdot$ repository-based measurability\\$\cdot$ activity \& maintenance (18-month rule) $\cdot$ scope coherence};
\node[final, below=of filt] (fin) {28 selected packages\\[1pt]\footnotesize\mdseries(measured in Chapter~\ref{ch:measurement-results})};
\node[excl, right=10mm of filt] (ex) {\textbf{Excluded:}\\MATLAB/Simulink, Unity, Isaac~Sim, RaiSim (not OSS)\\MRDS, USARSim (inactive)\\GPMP2, CHOMP (scope)};
\draw[arr] (s1) -- (cons);
\draw[arr] (s2) -- (cons);
\draw[arr] (s3) -- (cons);
\draw[arr] (cons) -- (filt);
\draw[arr] (filt) -- (fin);
\draw[arr, dashed, draw=orange!60!black] (filt) -- (ex);
\end{tikzpicture}
\caption{Candidate selection pipeline. Packages discovered from three source types are consolidated and ranked by mention count, then reduced by the filtering criteria of Section~\ref{sec:filtering-criteria} to the 28 measured packages. Mention counts are source-appearance signals only (Section~\ref{sec:mention-count-ranking}); representative excluded tools and the reason for each exclusion are shown at right.}
\label{fig:candidate-selection}
\end{figure}

\begin{table}[p]
\centering
\caption{Final software set with primary roles.}
\label{tab:final-software-set}
\singlespacing
\begin{tabular}{@{}ll@{}}
\toprule
\textbf{Software} & \textbf{Primary Role} \\
\midrule
AirSim & Simulator (autonomous vehicles) \\
Bullet/PyBullet & Physics engine \\
CARLA & Simulator (autonomous driving) \\
CoppeliaSim (V-REP) & Simulator (general robotics) \\
DART & Dynamics library \\
Drake & Toolbox (dynamics, planning, control) \\
Flightmare & Simulator (quadrotor) \\
Gazebo & Simulator (general robotics) \\
GraspIt! & Manipulation (grasping) \\
Habitat & Simulator (embodied AI) \\
MORSE & Simulator (academic robotics) \\
MoveIt! & Manipulation (motion planning) \\
MRPT & Library (mobile robotics) \\
MuJoCo & Physics engine \\
Newton Dynamics & Physics engine \\
ODE (Open Dynamics Engine) & Physics engine \\
OMPL & Motion planning library \\
OpenRAVE & Planning framework \\
OpenRTM-aist & Middleware (RT components) \\
OROCOS & Control framework \\
PhysX (open-source SDK) & Physics engine \\
Pinocchio & Dynamics library \\
Project Chrono & Multi-physics simulation \\
ROS~2 & Middleware and SDK \\
Simbody & Multibody dynamics library \\
SOFA & Simulation framework \\
Webots & Simulator (general robotics) \\
YARP & Middleware (humanoid robotics) \\
\bottomrule
\end{tabular}
\end{table}

The role column is necessary because the packages are not direct substitutes. ROS~2, MoveIt!, and OMPL are boundary cases because they are not full simulators in the same sense as Gazebo or Webots; Chapter~\ref{ch:selected-software-commonality} provides the detailed category and commonality analysis used to interpret them.

\section{Measurement Procedure}
\label{sec:measurement-procedure}

The measurement procedure follows the Domain~X method for collecting structured evidence about software packages in a selected domain~\citep{Smith2021}. For each of the 28 selected packages, the study gathered source code, documentation, publications where applicable, repository metadata, issue-tracker evidence, release information, and installation evidence. The measurement record identifies the specific repository or repositories measured, since some robotics projects are distributed across multiple repositories.

\subsection{Measurement Template}
\label{subsec:measurement-template}

The primary measurement instrument is the measurement template from Appendix~B of \citet{Smith2021}; the version used in this study is reproduced in Appendix~\ref{appendix:measurement-template}. The template contains approximately 108 fields when supplementary entries (free-text ``additional comments'' fields after each section and the GitHub-style and GitStats/SCC repository metric panels) are counted alongside the substantive questions. The substantive questions are grouped into the following sections:

\begin{enumerate}
    \item \textbf{Summary Information} (17 questions): software name, URL, affiliation, purpose, number of developers, funding, initial release date, last commit date, status, license, platforms, software category, development model, publications, source code URL, programming languages, and evidence of performance consideration.
    \item \textbf{Installability} (14 questions): installation instructions (presence, single-document organization, linear ordering, completeness, and assumed pre-installed dependencies), compatible operating system versions, installation automation, error handling, installation validation procedure, number of installation steps, operating system used during measurement, dependencies (presence of guidance and required version listing), and uninstall behaviour (problems and residue).
    \item \textbf{Correctness and Verifiability} (8 questions): requirements or theory documentation, correctness techniques, tutorial availability and quality, unit tests, and continuous integration evidence.
    \item \textbf{Surface Reliability} (5 questions): installation and tutorial breakage, error messages, and recoverability.
    \item \textbf{Surface Robustness} (2 questions): handling of unexpected input and edge cases.
    \item \textbf{Surface Usability} (4 questions): tutorial, user manual, user characteristics, support model.
    \item \textbf{Maintainability} (7 questions): version number, contribution guidance, available artifacts, issue tracking, issue closure rate, comment percentage, and version control.
    \item \textbf{Reusability} (2 questions): number of code files and API documentation.
    \item \textbf{Surface Understandability} (8 questions): code formatting, coding standards, identifier quality, constant usage, comment clarity, algorithm references, parameter ordering, and modularization.
    \item \textbf{Visibility and Transparency} (4 questions): development process definition, process documentation, development environment documentation, and release notes.
\end{enumerate}

Applying the same template to every package ensures that comparisons are based on a consistent set of observations rather than ad hoc impressions. Each quality section concludes with an overall impression score on a scale of 1 to 10, assigned using the scoring guidelines of the impression calculator in Appendix~C of \citet{Smith2021}. The impression scores summarize more than the recorded checklist answers alone: they also reflect observations made during measurement, such as the quality and currency of documentation, the clarity of error messages, and the overall installation and tutorial experience, which the mostly binary checklist fields do not capture. Two packages with identical checklist answers can therefore receive different impression scores; Chapter~\ref{ch:threats-to-validity} discusses the subjectivity this introduces.

\subsection{Repository Metrics}
\label{subsec:repository-metrics}

Repository metrics are collected separately from the quality impression scores. Three categories of repository data are gathered:

\begin{enumerate}
    \item \textbf{GitHub-style metrics:} stars, forks, watchers, open pull requests, closed pull requests, number of developers, open issues, closed issues, initial release date, and last commit date.
    \item \textbf{GitStats-style metrics:} number of text-based files, number of binary files, total lines, lines added, lines deleted, total commits, commits by year (last 5 years), and commits by month (last 12 months).
    \item \textbf{SCC-style metrics:} number of text-based files, total lines, code lines, comment lines, and blank lines.
\end{enumerate}

From these raw data, three processed measures are calculated:

\begin{enumerate}
    \item \textbf{Alive/dead status:} alive is defined as the presence of commits or version releases within the last 18 months from the measurement date.
    \item \textbf{Percentage of identified issues that are closed:} computed as closed issues divided by total identified issues.
    \item \textbf{Percentage of code that is comments:} computed as comment lines divided by total lines.
\end{enumerate}

\subsection{Measurement Execution}
\label{subsec:measurement-execution}

The Domain~X methodology recommends performing installation-based measurements in a controlled environment, typically a virtual machine, to ensure a clean and reproducible starting state~\citep{Smith2021}. For this report, installation measurements were performed on a Linux-based virtual machine. The operating system, installation steps, dependencies, and any issues encountered were recorded for each package.

The methodology also recommends a time budget to keep the overall measurement workload feasible. The full assessment is estimated at roughly 173 person-hours per domain, against a stated goal of about one person-month (160 hours), with 1 to 4 hours allocated per package for template completion and up to 2 hours for installation~\citep{Smith2021}. Measurement was conducted during May 2026. All time-sensitive data (last commit dates, issue counts, stars, forks, etc.) should be understood as snapshots taken during this period.

Several measurement questions require explicit interpretation rules:

\begin{itemize}
    \item \textbf{Last commit date} is recorded at the time of measurement and may change for active repositories.
    \item \textbf{Percentage of identified issues that are closed} is computed consistently as closed issues divided by total identified issues, using GitHub issue data where available.
    \item \textbf{Evidence that performance was considered} is assessed by searching software artifacts for explicit references to speed, storage, throughput, performance optimization, parallelism, multi-core processing, or similar considerations.
    \item \textbf{Percentage of code that is comments} is computed using the SCC tool's line-based method, as comment lines divided by total lines.
\end{itemize}

\subsection{Ranking with AHP}
\label{subsec:ahp-ranking}

After the quality measures are collected, the Domain~X methodology uses the Analytic Hierarchy Process (AHP) to support ranking across the nine quality categories~\citep{Saaty1980}. AHP is a decision-making technique that uses pairwise comparisons between alternatives to calculate relative scores. It organizes multiple criteria in a hierarchical structure and enables the combination of individual quality rankings into an overall ranking based on the relative priorities assigned to each quality.

In this report, AHP would be applied to the nine quality impression scores for the 28 software packages and would produce a ranking that depends on the chosen criterion weights. Sensitivity analysis would then assess whether the ranking is stable under reasonable variations in those weights. An elicited-weight AHP ranking is not computed here, because a defensible criterion-weighting elicitation has not been completed at the time of writing; Chapter~\ref{ch:measurement-results} therefore reports the per-package quality-mean ranking and a comparison with community indicators, together with an illustrative equal-weight AHP computation that demonstrates the weighting step.

\section{Commonality Analysis Procedure}
\label{sec:commonality-analysis-procedure}

The commonality analysis identifies shared and variable capabilities across the selected software packages, treating them as members of a software family as defined by \citet{Parnas1976} and \citet{Weiss1998}. The analysis supplements repository-based metrics by describing what the packages do and how their roles differ, providing context for interpreting the measurement results reported in Chapter~\ref{ch:measurement-results}.

Following the Domain~X methodology, the analysis proceeds in six steps:

\begin{enumerate}
    \item \textbf{Write an introduction} to the domain and the software family.
    \item \textbf{Write an overview of domain knowledge,} summarizing the key concepts and terminology relevant to MPS software.
    \item \textbf{List commonalities:} identify goals, assumptions, functions, artifacts, and domain concepts that recur across multiple family members. For example, many of the selected packages share the goal of simulating robot dynamics, the assumption that robot models are described using standard formats (such as URDF or SDF), and the function of providing an API for programmatic control.
    \item \textbf{List variabilities:} identify dimensions on which family members differ. These include software role (simulator, physics engine, planning library, middleware), supported platforms, licensing, programming language interfaces, documentation style, and intended use cases.
    \item \textbf{List parameters of variation:} for each variability, describe the possible values. For example, the software role parameter can take values such as general-purpose simulator, domain-specific simulator, physics engine, dynamics library, motion planning library, manipulation framework, or middleware.
    \item \textbf{Add terminology, definitions, and acronyms} that are standard in the domain, to ensure that the analysis is accessible to readers outside robotics.
\end{enumerate}

The procedure first groups the selected packages into broad categories based on their primary roles: robotics simulators, physics and dynamics engines, and planning or middleware tools. It then identifies capabilities that appear across multiple categories, such as support for robot model representation, simulated environments, motion planning workflows, dynamics computation, and integration with robotics ecosystems.

Because the selected packages span multiple roles, the commonality analysis does not force all tools into a single comparison template. Its purpose is to describe shared concerns and meaningful differences, not to claim that every package provides identical functionality. The analysis is based on checked evidence from documentation, repositories, publications, and other reliable project materials. Capabilities that could not be verified against these materials are not reported.

\section{Interview Component}
\label{sec:interview-component}

The Domain~X methodology includes a developer survey component designed to capture information that is not visible in public repositories: the development process as experienced by practitioners, the pain points they encounter, their perceptions of software quality threats, and the strategies they use to address those threats~\citep{Smith2021}. The methodology prescribes semi-structured interviews based on a list of 20 questions covering developer background, project organization, user understanding, development obstacles (past and present), documentation practices, and specific software qualities including maintainability, understandability, usability, and reproducibility.

For this report, the interview component requires research ethics approval from the McMaster Research Ethics Board (MREB), with the application submitted through the MacREM online system, before recruitment or data collection. The planned interview participants are software developers, maintainers, or programmers associated with the studied packages. Approved ethics materials and the interview protocol will be archived in Appendix~\ref{appendix:interview-materials}.

Chapter~\ref{ch:developer-interviews} records the planned reporting structure for the interview component. Until ethics approval, recruitment, and analysis are complete, that chapter does not report developer findings; it identifies RQ6--RQ9 as questions that require interview evidence.

The methodology recommends sending interview requests to developers from all packages in the study set to avoid selection bias. Prior applications of the methodology report response rates roughly in the 15\%--30\% range~\citep{Smith2021,Smith2024MI}; the exact rate for this report will depend on outreach timing and the response practices of the contacted developers. Interview contacts are typically identified through project websites, code repositories, publications, or institutional biography pages.
